\chapter{Discussion} \label{Discussion}

The results presented in Chapter \ref{results} are discussed in the sections that follow. First, the focus is on the three GFL strategies tested in the present study. These are compared in terms of cognitive processes, i.e. translation and PE times, screen activity tracking, and personal impressions of the participants. Product data, i.e. strategy applications and the success of their integration in the translation and PE process, are also discussed. The strengths and limitations of each strategy are described. Second, translation and PE are compared. These sections build the basis for discussing the implications of the study for both language professionals and MT research. Based on the results, a tentative decision tree for integrating GFL into the translation and PE process is proposed. Finally, a section is dedicated to the research process, describing the difficulties faced during data collection and analysis as well as the limitations of the present work. 

\section{GFL approaches compared} \label{sec:strenghts/limitations}

Translation and PE times increased when using neosystems. While time is surely needed to understand how the selected system works and also to memorise new pronouns, suffixes, and declension systems, no statistically significant difference was found in translation and PE times when comparing neosystems with the other approaches. However, this assignment took the longest for half of the participants, and the median translation and PE times were higher than in the other assignments. It is unclear whether a larger sample could yield different results and it must be highlighted that great differences in translation and PE speed were found among participants.

Gender-neutral rewording and gender-inclusive characters are rather common in written language and almost every participant reported regularly using one or both. As language experts, participants can adapt language and texts according to different registers, contexts, client specifications, and so on. This might be why the increased effort in the use of neosystems, e.g. the high number of searches and changes as well as the perceived cognitive load, does not correlate with significantly longer translation and PE times.

The analysis of the screen activity tracking elicited similar results to the analysis of translation and PE times. No difference was found between the first two assignments and the screen activity increased in the third. The difference was statistically significant in this case. This may be because neosystems are not common, and translators needed to consult various sources to be able to use them. The number of changes was rather high in each assignment with almost no differences among gender-inclusive characters and neosystems. Currently, there is no standard for the use of gender-inclusive characters. As a matter of fact, participants used different characters and/or combined them differently. When in doubt, they also consulted resources, which sometimes led to changes, e.g. from ``der*die'' (masculine*feminine article) to ``die*der'' (feminine*masculine article). When applying neosystems, participants were for instance not sure how to use pronouns. These challenges led to a high number of text modifications both in the second and third assignments.

Screen recordings also provided interesting insights into the participants' general approaches to translation and PE. The majority searched online for information concerning the TV series and the non-binary actors/characters mentioned in the study texts. Some even looked for photos. While more information may be useful to produce a better translation, searching for a person's photo online is not likely to assist in the process. As explained in Section \ref{sec:sex/gender}, there is a distinction between gender identity and gender expression, hence using the physical appearance in photos or in general to infer a person's pronouns may lead to misgendering and identity invalidation. Participants also highlighted gender references to be adjusted. They often conceived the translation and PE process as a separate step from the use of GFL. Some participants in the PE task adjusted gender references first and then worked on the whole text. This seems to suggest that the use of GFL has a partially negative impact on cognitive processes at least at this stage of familiarity with GFL.


Increased translation times do not correlate with higher screen activity tracking. This may also be because the participants' speed and productivity greatly differed as demonstrated by the high standard deviation in translation and PE times. It seems that some participants are faster than others, even though they make several changes or consult numerous resources. This may also be due to different levels of work experience (spanning from 3+ to 20+ years) and/or greater knowledge of GFL. In the pre-study questionnaire, participants were asked whether they had previous knowledge of GFL and participants could opt for ``yes'', ``little'', and ``no''. While every participant in the study had at least a little experience, no more fine-grained assessment of their knowledge was carried out. During the interviews, it clearly emerged that some participants had greater awareness of GFL use and issues. In future research, an assessment of GFL ``skills'' could reveal whether they have an impact on translation and PE times, effort, and target text quality. 

Interviews suggest a preference for gender-inclusive characters, which is considered a relatively straightforward approach. Nevertheless, participants commented that the frequent use of characters may impact the target text readability. Opinions concerning gender-neutral rewording were more contrasting because finding gender-neutral alternatives to commonly used terms is seen both as a strength and a limitation of this approach. One can use the existing language inventory without introducing disruptive forms but this may also lead to the use of very uncommon terms or sentence structures. This may be the reason for the higher number of pauses and the lower number of searches in the first assignment. Participants commented that the use of neosystems was cognitively demanding, comparing them to a foreign language whose functioning is rather unclear. Their attitudes towards such systems differed: some were particularly negative, e.g. one described them as ``raping the German language'', whereas others were generally positive even though participants argued that they would need further training to understand and apply them correctly. 

Finally, the target text annotations could be used for triangulation with the process data. Many strategies are possible in the first assignment, e.g. gender-neutral words, passive constructions, and repetitions of proper names. For instance, participants used several neutral forms to translate the same gender-specific terms, e.g. ``Schauspieltalent'' (acting talent), ``Schauspielstar'' (acting star), ``Schauspielprofi'' (acting professional) to avoid using gender-specific ``Schauspieler'' (male actor) and ``Schauspielerin'' (female actor). In addition, they reworded whole text passages, often changing perspectives, abstracting, adding, or omitting information. This reflection and revision of one's work is among the core competencies of translation professionals and differentiates them from MT systems, i.e. the creativity and the capacity to adapt texts to different cultures, requirements, and specifications. 

Participants showed a clear preference for gender star (*) in the second text, although its use was inconsistent among them. During the study, several participants stated that they were unsure about the order of masculine and feminine forms, which may also have impacted the number of changes that were made, in part, to alter the order of gendered forms in articles and pronouns. The translation analysis also shows that sometimes gender star was combined with another character, i.e. slash. This could be a result of changes in language use. Once common gender-inclusive forms included Binnen-I, e.g. \textit{LeserIn} (masculineFeminine noun), and slash, e.g. \textit{Leser/in} (masculine/feminine noun). These forms are now considered outdated because they only include the male and female genders and are being replaced by other characters. Furthermore, there is no norm concerning their use. This is also an aspect that was discussed in the interviews, with participants expressing the desire and need for rules to ensure systematic use of characters. Finally, the practical implementation of some forms may be difficult, mostly when the word stem changes after adding the feminine suffix to the masculine form, e.g. ``Arzt/Ärztin'' as in the translation experiment. 

The most used neosystems were the Sylvain system and the Ens-forms. The former was selected due to its completeness, the second due to its ease of use. All participants made mistakes in their application, which seems to confirm the increased effort in comparison to the first two assignments. The consensus that this strategy is cognitively challenging was suggested by the target text analysis as well. The number of mistakes made in the application of neosystems was very high, whereas there were fewer, if any, application errors for the strategies in the first two assignments.

The analysis of the translation and PE times as well as screen activity tracking elicited interesting results, but interviews provided an opportunity to discuss the translation and PE process in more detail. They offered valuable insights that help describe the strengths and limitations of the strategies investigated in the present work. These are briefly summarised below.

Gender-neutral rewording requires rethinking the translation of very simple terms. This may be regarded as what translation practice is about: rewriting (\cite{lefevere1992translation}) or ``dire quasi la stessa cosa'' (say almost the same thing, \cite{eco2016dire}). However, constantly avoiding gender-specific terms leads to translation choices that are not always automatic and/or natural-sounding, as argued by the participants who were not satisfied with their translations. Participants stated that target texts may sometimes sound impersonal -- e.g. because of the use of passive constructions -- and/or may shift meanings. In terms of the translation and PE process, these may be the reasons why they made more pauses in the first text than in the other assignments during the experiment. A relatively high occurrence of changes was also observed, although to a lesser extent than in the other assignments. Furthermore, repetitions of proper names must be repeated to avoid pronouns, which led participants to rewrite whole text passages. 

%Nevertheless, texts containing neutral formulations are probably the most comprehensible and may be perceived as less ideological since gender is not a particularly salient feature of texts containing such formulations. Further research on this strategy should elicit results on whether and to what extent gender-neutral forms go unnoticed in different text typologies, e.g. administrative/legal texts. 

%Finally, since texts written utilising this approach may sound more abstract and/or impersonal, this might not be a desirable solution if the purpose of the source text is to promote gender diversity or to depict non-binary individuals. In this case, other approaches such as gender-inclusive characters and/or neosystems might be more appropriate.

Gender-inclusive characters, especially gender star, are the most common GFL strategy in German-speaking countries. Their use is relatively straightforward since usual masculine and feminine forms are concatenated with a character. However, there is no established norm to do so. Even though participants showed a strong preference for gender-inclusive characters, they would definitely mix it with rewording to avoid a strong negative impact on readability. Indeed, these two approaches were often combined in the second assignment. This is due to considerations of readability, but it also has to do with the nature of translation where a one-to-one correspondence between source and target language is not always possible and/or desirable. The relatively straightforward combination of masculine and feminine forms may be a limitation of gender-inclusive characters. Results from the interviews confirm that participants explicitly think of masculine and feminine forms of words to combine them with a character. This mindset reinforces the gender binary, while new gender-fair forms may be more inclusive of all genders. Using gender-inclusive characters could thus be just a first step towards more GFL because they are already common to some extent and may therefore be met with less resistance by a broad public. 

While neosystems may be the best solution to foster the visibility of non-binary genders because they are new, unseen forms specifically devised for this, they are largely unknown. Only one participant in the experiment used a neosystem in their everyday life. Others had never or hardly heard of them before. In addition, neosystems also interfere with classical grammatical rules, so the number of searches was extremely high during the translation and PE of the third assignment. Participants had to repeatedly search for instructions regarding their use and then ensure the correctness of their translation solutions. Limited knowledge and difficulties/uncertainties also affected the number of changes, which was the highest for all assignments. 



\section{Translation and PE compared}
PE was faster than translation in the first and second assignments. This is not surprising since a great body of research on PE already confirms this (e.g. \cite{jia2019does}). However, no statistically significant difference between the two tasks was found in the third assignment. This also seems to suggest that the use of neosystems involves higher effort, thus reducing the advantages in speed brought about by the automated translation. Participants in the PE group felt that MT increased their productivity. According to them, this was due to the good quality of the machine translations despite their biased nature.

Differences in screen activity tracking were less clear-cut. First, they were not fully comparable. As already explained in Chapter \ref{results}, the MT outputs contained masculine generics, misgendering, and literal translations of singular \textit{they}. This meant that participants in the PE group had to revise nearly all gender references, something that the colleagues in the translation group did not have to do. Therefore, in a further analysis, the number of changes was excluded from the total count of the activities performed. Median values seem to suggest higher activity for translation in the first assignment, no differences in the second, and higher activity for PE in the third. Such differences among tasks, however, were not statistically significant.

The retrospective interviews also showed few differences between the tasks. Participants in the PE group showed slightly more negative attitudes towards neosystems. The real difference concerned perceptions of the utility of MT in the assignments. Translators did not generally use MT and would not consider it for texts similar to those in the study. Post-editors felt that, despite the effort to correct biased outputs, they were more productive.

The most interesting differences between tasks emerge in the target text analysis. In the first and second assignment, translators made more mistakes. The most common mistake in the first was using the wrong term ``Studierende'' to refer to high-school students. In the second, some participants interpreted singular \textit{they} as plural and also made mistakes when translating the challenging term ``doctor''. Translators reworded more often than post-editors in the second assignment, paying attention to readability and target text fluency. This may be due to the influence of the given MT, which constrains creativity and the ability to extensively adapt texts -- a common PE issue mentioned by the study participants. Finally, participants from the translation group selected the Ens-forms since they were considered the easiest neosystem. Participants from the PE group selected the Sylvain system because it is the most developed and complete. This difference in the selection may be explained by the fact that post-editors have more capacity to concentrate on gender since they already have a good translation output that must simply be revised and stylistically adapted. Translators have to produce the target text without the support of MT, thus selecting the neosystem which presumably takes less time and effort.

\section{Implications for translation practice and theory}
%adattare testi sulla base di molti fattori diversi --> qualcosa che viene gia fatto anche se non in termini di linguaggio di genere 
The results of this study have several implications for both translation practice and theory. The first part of this section focuses on the implications for language professionals, including the context dependency of translation choices, the specialist knowledge required for producing gender-fair translations, the constant training needed to keep up with linguistic change, and the strengths and limitations of MT. The second part focuses on translation theory, exemplifying how this work expands the research object of feminist and queer translation and reinforces the conceptualisation of translation as a process of cross-linguistic, political, ideological, and ethical negotiation. 

During the post-study interviews, participants discussed challenges concerning the use of GFL. They stated that translation choices, especially with regard to GFL, are extremely context-dependent as also observed in joint studies (\cite{burtscher-2022-genderfairmt,gromann2023participatory}). Therefore, several factors must be taken into account when translating. These include: 
\begin{itemize}
\setlength\itemsep{-0.5em}
    \item \textbf{The source and target culture}: what are common views on gender? Do they differ from the source to the target culture?
    \item \textbf{The source and target language}: how is gender expressed in the source and target language? Are there differences in grammatical structures? Are there established gender-fair forms? How do these differ from source to target? What are the implications of the use of a specific strategy in contrast with the source language (and culture)? 
    \item \textbf{The source and target text}: what is the purpose of the text? Does gender play an important role? Must it be visible? Is the purpose of the target text the same as that of the source?
    \item \textbf{The source and the target audience}: are all genders addressed? Are specific genders addressed? Does or should the text address a broad audience?
    \item \textbf{Accessibility} (connected to the audience): is the strategy comprehensible? To what extent does it affect readability? Can screen-readers pronounce it? Can it be pronounced at all or does it work in the written language only?
    \item \textbf{Referents}: which people are explicitly mentioned in the text? Are they real people or fictional characters? Is it possible to gather information on their pronouns? Is it possible to contact them and ask which pronouns to use in the translation? 
    \item \textbf{Clients}: what are their attitudes toward GFL? Do they have preferences for a specific strategy? Are there any strategies that they absolutely reject?
    \item \textbf{Time constraints}: has the translator time to focus on the gender aspect? Are they paid for research, e.g. looking for pronouns used by people mentioned in texts and contacting them to discuss a translation strategy?
    \item \textbf{Resources}: which materials, e.g. instructions regarding GFL, guidelines, translation memories, etc., are available to translators? Do they have contact with non-binary and queer community members to make informed decisions?
    \item \textbf{GFL approach}: what are the advantages and disadvantages of each approach? Which strategies are more common in reference to a non-binary individual? Which strategies impact readability? Which strategies may be less comprehensible? And which are met with greater resistance?
\end{itemize}
These factors complicate the use of GFL in translation and highlight the impossibility of a one-size-fits-all approach. This is one of the main implications of this study for language professionals. Decisions must take account of all the above aspects. Hence, different approaches and/or strategies may be appropriate since they all have strengths and limitations as described in Section \ref{sec:strenghts/limitations} and acknowledged by participants during the post-study interviews.

GFL may be regarded as a form of subject-matter expertise due to the large variety of available strategies and their constantly evolving nature as well as the terminology related to gender/sex identities. Such knowledge requires constant training, possibly establishing contacts with interested communities as put forward by \textcite{attig2023call}. Language evolves rapidly and the acceptability of strategies changes. Those that were once common may now be outdated, e.g. the aforementioned Binnen-I, whereas new strategies have been and are being developed. Furthermore, the use of terminology is shifting. Terms such as ``transexual'' are pathologising and nowadays considered unacceptable (\cite{lopez2019tu}). Hence, based on the results of this study and the knowledge gained by the author while working on the present monograph, a tentative best practice model is proposed in the next section. However, more research on this topic is needed and GFL use might change rapidly.

Finally, a major implication concerns the use of technology. While translators' attitudes towards MT are still quite negative, they improve with the use of technology. PE can be a valuable tool in improving translators' productivity and, in part, ensuring better translation quality even though the MT outputs are severely biased. Conversely, the value of human translation is demonstrated by bias in MT, and by the translators' skills in creatively avoiding gendered elements as well as adopting different strategies based on clients' specifications. Translators can rapidly adapt to linguistic change and even foster it through their work, possibly becoming agents for social change, not only through their translations but also their contacts with clients. For example, they may explain why GFL should be used on websites, in product specifications, and in other text types, showing how inclusion could be important for the client's business. However, this should not be the only reason for adopting GFL. Furthermore, translators can understand identity terms, e.g. non-binary and transgender, and preserve them in their translations in comparison to MT. They can also produce many different versions of the same texts, adapting them to different requirements including different GFL approaches. To date, this is something MT systems can only accommodate to a very limited extent, if at all (whereas LLMs show promising results).

%-------------NEW--------------------
At a broader level within the discipline and translation theory, this monograph expands the scope of feminist and queer translation research for two reasons. First, it focuses on a largely neglected and marginalised group, namely non-binary individuals. Second, it broadens the analysis of feminist and queer translation strategies to include fields such as audiovisual and news translation, which have received less attention in the literature (\cite{baer2017introduction,von2020routledge}). Concurrently, the study's results reinforce the conceptualisation of translation as a process involving cultural, political, and ethical considerations.

The linguistic representation of non-binary individuals has been underexplored in the field of queer translation. Queer translation has traditionally concentrated on sexual orientation, as exemplified by \textcite{harvey1998translating}, who examined the translation of camp talk -- a sociolect specific to gay men. It also has addressed gender identity representation, but primarily from a binary perspective. For example, \textcite{rose2016keeping} investigates how to preserve the gender identity of transgender individuals in the translation of French memoirs into English, where both masculine and feminine gender markers are used for the same transgender person.

The analysis of non-binary GFL in this study, through its translation and PE processes, reveals the necessity of breaking traditional grammatical rules and experimenting with language. This underscores how language and translation can function as queer practices, deconstructing and reconfiguring linguistic boundaries by introducing new linguistic elements, such as neopronouns and neomorphemes, to challenge existing binary gender norms and assumptions. As discussed in the interviews, translators must negotiate various co-textual and contextual factors, including the shifts in meaning induced by GFL usage, target audience considerations, client specifications, and personal and societal views.

Given its controversial and politically/ideologically charged nature, GFL can face harsh criticism and resistance (\cite{vergoossen2020four}). Additionally, the unfamiliarity with neopronouns and neomorphemes can impact the readability and comprehensibility of texts. The translation process thus becomes more visible due to the lack of equivalence between language structures. While English is generally neutral, German-speaking language professionals must choose a GFL strategy to avoid producing a misgendering translation.

Feminist and queer translation research has traditionally focused more on literary works, with audiovisual products receiving less attention, and news translation being even more neglected. The concept of social responsibility has been explored in translation and interpreting within sensitive contexts such as medicine, law, and migration (\cite{drugan2017translation}). This study illustrates how gender is a highly sensitive area too, where translators bear the responsibility for a respectful representation of marginalised groups. The increased visibility of non-binary individuals is leading to their representation in popular audiovisual and news media, which necessitates translation across numerous languages. This can potentially foster social change by normalising the notion that gender is not binary and by reducing gender discrimination.


%-------------END NEW--------------------

\section{A tentative best practice model for GFL in translation and PE}
Based on the previous considerations and knowledge gathered by the author while writing this monograph, a simplified and tentative decision tree for selecting the most appropriate GFL strategy is proposed in Figure \ref{fig:taxonomy}. This is by all means an initial proposal that could and should be refined, especially after discussing it with members of the queer and non-binary community and language professionals. The decision tree is discussed below in the context of translations from English into German, thus further commenting on the implications of the results from this monograph.


\begin{figure}
\footnotesize
\caption{Decision tree for the selection of the appropriate GFL strategy}
\label{fig:taxonomy}
\begin{forest}for tree={node options={align=left},draw},forked edges
  [{Does the text mention a queer/non-binary person?}, text width=7cm,align=center
    [{Is there an\\
    equivalent\\
    term for their\\
    gender in your\\
    language/culture?},tier=disruptive,text width=2.5cm,align=left, edge label={node[midway,xshift=-3mm]{\textbf{yes\strut}}}
        [{Explain the difference\\
        between the source\\
        and the target language.\\
        Look for resources, e.g. \\
        queer blogs and \\
        language guidelines,\\
        to make an informed\\
        decision. Discuss\\
        your decision with the\\
        person.}, text width=3.2cm, align=left,tier=explain,fill=black!10!white, edge label={node[midway,xshift=-3mm]{\textbf{yes\strut}}}]
        [{Is the person real?}, edge label={node[midway,xshift=3mm]{\textbf{no\strut}}}
            [{Are there established\\
            terms to refer  to them?},align=left, edge label={node[midway,xshift=-3mm]{\textbf{no\strut}}}
                [{Provide an\\
                explanation\\
                after consulting\\
                resources, e.g.\\
                queer blogs.},text width=2.1cm,align=left,fill=black!10!white, edge label={node[midway,xshift=-3mm]{\textbf{no\strut}}}]
                [{Use these}, tier=use,fill=black!10!white, edge label={node[midway,xshift=3mm]{\textbf{yes\strut}}}]
            ]
            [{Can you find information\\
              about and/or contact  them?}, text width=3.7cm,align=left,thick, edge label={node[midway,xshift=3mm]{\textbf{yes\strut}}}
                [~,edge label={node[midway,xshift=-3mm]{\textbf{no\strut}}},name=source, draw=none,text width=.5cm]
                [{Are there equivalent\\
                strategies in your\\
                language to refer
                \\to them?
                \\~},text width=2.7cm, align=left, edge label={node[midway,xshift=3mm]{\textbf{yes\strut}}}
                    [{Explain the difference between the source\\
                    and the target language. Look for resources,\\
                    e.g. queer blogs and language guidelines,\\
                    to make an informed decision. Discuss your\\
                    decision with the person.}, text width=6cm,align=left,tier=explain,fill=black!10!white, edge label={node[midway,xshift=-3mm]{\textbf{no\strut}}},name=target]
                    [{Use these},fill=black!10!white, edge label={node[midway,xshift=3mm]{\textbf{yes\strut}}}]
                ]
            ]
        ]
    ]
    [{Does gender play an important role?}, edge label={node[midway,xshift=3mm]{\textbf{no\strut}}}
        [{Reword sentences in\\
        order to avoid gender.},fill=black!10!white, text width=3cm, edge label={node[midway,xshift=-3mm]{\textbf{no\strut}}}]
        [{Is the target audience \\
        wide or does the text\\
        address queer people only?}, text width=4cm,align=left, edge label={node[midway,xshift=3mm]{\textbf{yes\strut}}}
            [{Opt for a mix of   gender-neutral\\
            and             (common) inclusive \\
            strategies.  Motivate   your \\
            selection.}, text width=4.5cm,align=left,fill=black!10!white, edge label={node[midway,xshift=-4mm]{\textbf{wide\strut}}}]
            [{Use disruptive\\
            strategies,  e.g.\\
            neopronouns,\\
            neomorphemes,\\
            or characters,\\
            if available},tier=disruptive, text width=2.2cm,align=left,fill=black!10!white, edge label={node[midway,xshift=5mm]{\textbf{queer\strut}}}]
        ]
    ]
  ]
  \draw(source.north)--(target);
\end{forest}

% \begin{tikzpicture}
%   \node(start) [draw, diamond, aspect=1,text width=2cm,font = {\footnotesize},align=center] {Does the  text mention a queer{\slash}non-binary person?};
%   \node(equivalent)[right=of start,draw, diamond, aspect=1,text width=2cm,font = {\footnotesize},align=center]{Is there an equivalent term for their gender in your language/culture?};
%   \node(real)
% \end{tikzpicture}
%
% NO
% Does the text mention a queer/non-binary
% person?
% YES
% Does gender play an
% important role?
% NO
% YES
% Is the target audience
% wide or does the text
% address queer people
% only?
% NO
% Provide an explanation
% after consulting
% resources, e.g. queer
% blogs.
% Reword sentences in
% order to avoid gender.
% Are there established
% terms to refer to them?
% NO
% Is there an equivalent
% term for their gender in
% your language/culture?
% YES
% Ensure that the term
% used is current and not
% pathologising
% Is the person real?
% NO
% Explain the difference
% between the source and the
% target language. Look for
% resources, e.g. queer blogs
% and language guidelines, to
% make an informed decision.
% YES
% Usethese
% YES
% Can you find
% information about and/or
% contact them?
% YES
% WIDE
% Opt for a mix of gender-
% neutral and (common)
% inclusive strategies.
% Motivate your selection.
% QUEER
% Use disruptive strategies,
% e.g. neopronouns,
% neomorphemes, or
% characters, if available
% Are there equivalent
% strategies in your
% language to refer to them?
% NO
%  YES
% Use these
% Explain the difference
% between the source and the
% target language. Look for
% resources, e.g. queer blogs
% and language guidelines, to
% make an informed decision.
% Discuss your decision with
% the person.
\end{figure}


An important criterion when deciding if and how to use GFL is whether non-binary individuals are referenced. While in the present work it is argued that masculine generics represent a sexist use of language and should be avoided, GFL must actively be used in reference to non-binary people. This is not a question of personal preference and/or ideology. Intentionally assigning the wrong linguistic gender to a person is an act of identity invalidation and misgendering that can have serious negative impacts on the mental health of trans and non-binary people (\cite{mclemore2018minority}). 

When a non-binary person is not directly mentioned in the source text, one should consider whether gender plays an important role in the target text. If not, for instance in product instructions, translators could avoid masculine generics by using gender-neutral rewording. If the gender element is important in a text, the strategy selection is dependent, amongst others, on the target audience. In an administrative text, addressing a wide audience, gender-neutral rewording ensures that the text is highly comprehensible and readable. The use of this strategy may also go unnoticed and hence reduce or eliminate perceived ideology in the text.

Texts written using this approach may sound more abstract and/or impersonal, which might not be a desirable solution if the purpose of the source text is to promote gender diversity or to depict non-binary individuals. In this case, other approaches such as gender-inclusive characters and/or neosystems might be more appropriate. Translators could explain the rationale for the strategy selection as a form of awareness raising. Gender star was the most selected gender-inclusive character strategy in the present work, so its use is highly recommended. While combined forms, e.g. ``die*r'' are a viable solution, these are not always possible, e.g. in the case of possessive pronouns ``sein*ihr'' (his*hers). If consistency is desired, masculine forms should always be used first, e.g. ``der*die Schauspieler*in'', (masculine*feminine article + masculine*feminine noun). Otherwise, the use of feminine articles and pronouns first is recommended as a strategy against androcentrism in language. So, if the word stem changes when adding the feminine suffix ``-in'' as in ``Arzt/Ärztin'', the feminine word stem should be preferred as in ``Ärzt*in''. It is also important to stress that gender-inclusive characters should be sparsely used in order not to impact reading flow and readability. A mix of these two approaches was also indicated as the preferred solution by the participants in this study.

In texts written for a queer audience and/or for queer activism purposes, more disruptive strategies, such as neosystems, could be actively used to linguistically include all genders and to question the gender binary. There are many different neosystems, and the selection could be based on their completeness, e.g. the Sylvain system, or inherent ease of use, e.g. the Ens-forms. If possible, it would be best to decide after discussing with members of the queer/non-binary community since other systems may be gaining popularity and/or may be preferred. To the best of my knowledge, neosystems are extremely rare in German-speaking countries and non-binary individuals normally use different neopronouns, e.g. \textit{nin} or \textit{xier}, and gender-inclusive characters to talk and write about themselves. 

If a specific person is directly mentioned in the source text, the first important step is to use the correct terminology. The translated identity terms must be current and respectful. If this is a difficult decision, there are plenty of resources at hand, e.g. queer blogs or language guidelines such as \textcite{apagenderenglish}. If no equivalent terminology is found after consultation, an explanation can be provided together with a tentative translation. In German, English language terms are often borrowed.

If the non-binary person mentioned in the text is real, information on their pronouns could be searched in order to decide on one or more strategies. If no equivalents are found, a good solution would be to explain the difference between the source and the target language, e.g. explaining how the pronoun \textit{they} can be used in the singular to refer to a specific person whereas this is not possible in German, and select a strategy after consulting resources. If possible, the non-binary person mentioned in the text should also be contacted before making a decision. As already mentioned, in German-speaking countries non-binary individuals often use different neopronouns and gender-inclusive characters to talk and write about themselves. This could be a good solution. If neosystems are used, the translation should be proofread by at least one other person with good knowledge of non-binary GFL in order to avoid mistakes which, in the present study, were very common. 

In certain cases, like audiovisual translation and app localisation, there are character limits. Gender-neutral rewording and/or the use of gender-inclusive characters may not always be possible. To overcome these obstacles and avoid misgendering and/or uncommon linguistic strategies, language professionals must be creative. In addition, some markup languages might interpret an asterisk as a command to print the following text in italics. Further research is needed on GFL that takes into account multi-modality and the challenges of audiovisual translation or localisation.


\section{Implications for MT research}
The results of the present study have important implications for further research in MT debiasing and training. These include shifting away from the current one-to-one paradigm in which a single translation is produced for one source text, the need for personalisation, the creation of gender-fair datasets, and constant (re)training as well as the need for participatory research. 

Recently, the one-to-one paradigm of MT has been shifted to a one-to-many approach where two or more translations are produced for one single source text. For instance, Google Translate started to show the feminine and masculine translation for gender-ambiguous sentences in 2018 (\cite{googlebias}) whereas \textcite{deepl} allows users to select the tone (informal/formal) of the translation in some languages. This should be similarly implemented for GFL, thus showing different outputs, e.g. a gender-neutral and a gender-inclusive version of the target text.

In order to do so, however, MT systems should be able to recognise terms like singular \textit{they} as such when translating from English into other languages. As seen in this study, singular \textit{they} was always translated as the third-person plural pronoun ``sie'' in German, thus misrecognising the gender of the referent. Plug-ins that allow for context disambiguation have been introduced (e.g. \cite{mechura-2022-taxonomy}) but still do not work well for non-binary genders and very few studies tried debiasing beyond the binary (\cite{saunders2020neural}). The lack of gender-fair datasets that could be used during training remains a major problem. Currently, bilingual corpora used for MT training are often composed of speeches from the European Parliament, and analyses of these materials show that most of them are given by men, with only about a third by women (\cite{vanmassenhove2019getting}). Non-binary people are extremely underrepresented in the training data.

While rewriting models have been proposed to implement GFL in MT systems, these are currently limited to certain strategies and their outputs still need corrections (\cite{amrhein-etal-2023-exploiting}). Another problem of such systems is particular gender-specific terms as shown in the present study. Some nouns are difficult to translate because there is no established gender-neutral alternative or, as in the case of ``nanny'', they generally have the masculine or the feminine form only. Therefore, MTPE may at least for the moment be the best solution to produce gender-fair datasets which are inclusive of genders beyond the binary or entirely neutral regarding gender references. Finally, rewriting models reduce representation harms but not the allocative ones: women and non-binary people would still need to generate an automatic translation and then use such models to produce a gender-fair version of the MT outputs, thus making MT less accessible for them. 

Another major challenge for MT research is linguistic change. As already mentioned, conceptions of what constitutes GFL shift, and strategies or terms that once were common are no longer used. Now that non-binary people are more visible, many GFL strategies are being proposed and/or redefined. MT systems would thus need constant updates. The present design and training procedures of MT systems make updating existing models quite cost- and resource-intensive and have a considerable carbon footprint (\cite{shterionov2023ecological}). Future research should thus prioritise more efficient methods for training and updating an MT model. Besides, non-binary people use different pronouns and different strategies to talk about themselves. MT systems should allow for personalisation, enabling users to select their preferred strategy.

The above points also require a shift in research on gender (de)bias(ing) in MT. Participatory research approaches (\cite{burtscher-2022-genderfairmt,gromann2023participatory}) are needed to include communities most affected by the use of GFL or lack thereof, i.e. queer and non-binary individuals. When it comes to marginalised communities, their exclusion from decision-making processes could be regarded as a form of epistemic violence (\cite{chakravarti1988can}). Failing to include non-binary people in deciding which GFL strategies are implemented further silences them, depriving them from their right to linguistic self-determination. Also, language professionals who might have different insights into strategy feasibility and comprehensibility should be involved in research to ensure a beneficial exchange between different stakeholders.

\section{Research process: Positionality and study limitations} \label{sec:limitations}

The research process involves constant decision-making from the topic selection and the research question to the data collection and analysis. I would like to first describe my positionality because it has an impact, both conscious and unconscious, on the decisions made regarding study design and data analysis. Finally, I will discuss the limitations of the present work. 

The use of GFL is often regarded as political and ideological. It also relates to one's beliefs concerning gender and language, mostly linguistic conservatism, e.g. stances on preserving the language. While some publications describe common arguments used against GFL (e.g. \cite{vergoossen2020four}), the objective of the present work was to investigate how language professionals apply it. GFL can be regarded as a language variety and as such it is of interest for language professionals. Besides, since the visibility of non-binary people is increasing, translators are and will be increasingly confronted with texts in which non-binary people are mentioned, e.g. news reports, TV series, and literary fiction. Hence, understanding the implications and the process of gender-fair translation is of utmost importance for TS and MT.

In the present work, there is no intention to discuss GFL in a politically and ideologically motivated manner. Nevertheless, it is a fact that refusing to respect a person's gender identity constitutes misgendering and this negatively affects the mental health of transgender people. If misgendering is a form of discrimination and epistemic violence, GFL use is a form of respect. I would argue that refusing to use GFL and intentionally misgendering a person is actually more ideological than the contrary. Language is an instrument to actively construct, negotiate, and represent human identity and social reality, and intentional misgendering aims to invalidate a person's identity.

This view of language reflects my positionality as a researcher and as such my work is politically and ideologically motivated. I refuse essentialist views of gender as innate and biologically determined. I embrace the constructionist view, recognising that gender as well as research are historical, cultural, and social constructs which nevertheless have an impact on how people are perceived and in turn may (or may not) experience forms of discrimination. I also recognise how language is the result of social conventions and how it (re)produces social norms and assumptions. These norms and assumptions, e.g. the gender binary, can also be challenged and changed through language, of course.

%I briefly discuss my privilege because I conduct research that involves and perhaps affects a marginalised group\footnote{The term marginalised is used to highlight that discrimination and exclusion are perpetrated because of unequal power distribution among people in society. It is preferred over minorities since discriminated groups, for instance, women are not a minority.} to which I do not belong. For this purpose, the concept of ``intersectionality'' needs to be introduced. The term was coined by \textcite{crenshaw1989demarginalizing} and describes how race, class, gender, and other individual characteristics overlap with one another. An example to understand how the concept works is the lawsuit by Emma DeGraffenreid against General Motors (GM) in 1976. GM was accused of discriminating against black people but won the lawsuit. When the judges verified that the company hired black people and white women, they believed there was no discrimination against black women, ignoring that black women face discrimination precisely because they are both black and women (\cite[142f]{crenshaw1989demarginalizing}). One analogy that can be used to understand intersectionality is that of crossroads. Every road represents a line of identity, for example, gender, sexuality, race, and class. The unique combination (or intersection) of these roads describes the multiple oppression faced by a single person. Fig. \ref{fig:intersectionality} offers an intuitive instrument to describe the different axes of one's identity. I am an Italian, white, able-bodied, cisgender, middle-class, researcher at a public institution in Austria. I never experienced forms of oppression and discrimination like people from marginalised groups. This also means that I have blind spots of which I am not always aware and which influence my beliefs and my research as well as the interpretation of my research results.

%\begin{figure}[ht]    
%\centering
%\includegraphics[width=0.8\textwidth]{figures/Intersectionality_wheel_of_power_and_privilege_from_University_of_British_Columbia.png}
  %  \caption{The wheel of power and privilege}
   % \label{fig:intersectionality}
%\end{figure}

The selected method comes with strengths and limitations. To the best of my knowledge, methods from TPR have not yet been applied to research on GFL. As discussed in \textcite{nettt2022}, most of the studies on GFL translation consist of analyses of translated audiovisual products in which non-binary individuals are featured. Such analyses focus on the use of GFL or lack thereof (\cite{lopez2022trans} and \cite{misiek2020misgendered}). In some cases, non-binary participants gave interviews and were asked which language strategies they use and/or would like to find in translations (e.g. \cite{vsincek2020ona}). Very little research addresses the translation process, generally involving students who compile questionnaires on the perceived difficulties of different translation tasks that require the use of GFL strategies (\cite{bailach2021como}). The selected methodology allows the collection of an interesting and especially multi-faceted dataset, proving this study to be among the few that provide in-depth results on this area of research.

As is common in TPR, data were collected and triangulated using a mixed-method approach. Data from the retrospective interviews, observational protocols and segment-based target text annotation were mainly textual and statistical analyses were performed to summarise the results. Screen recordings allowed translation times to be measured and screen activity tracking data to be analysed. The methodologies adopted allow the analysis of the collected data from different perspectives, ensuring a systematic and valid approach to empirical research.

%spiegare che le varaibili non sono state controllate perche importante e il realismo e inoltre cerano pochi partecipanti a disposizione. Finally, le interviste sono la parte piu preziosa dei dati e i partecipanti hanno dovuto tradurre proprio per poter avere insights in cio che serve loro. = mettere le persone della condizione di lavorare per ottenere il loro punto di vista.

%The setting in which the study was conducted was as realistic as possible. First, only translation professionals could participate in the study. Literature on the distinction between professional and non-professional translators is vast and it is beyond the scope of the present thesis to contribute to this discussion. Thus, in the context of this study translation professionals were defined as those with a relevant degree and/or at least two years of work experience in the translation industry. Second, participants could work in their office at home since they are freelancers. Third, they received realistic tasks for which they were also paid according to the recommendations by UNIVERSITAS Austria. Finally, they could choose whether to use or not Tools and online resources such as dictionaries and websites among others. 

The selected research design and methodology are also motivated by experimental realism and the sensitivity of the topic. A great variety of studies in the field of TPR adopt keylogging to analyse keyboard activity and eye-tracking to gain insights into the cognitive load. These methodologies are more precise and objective, but they require participants to be in a laboratory and hence impact ecological validity (\cite{o2009eye}). Study candidates are required to work in a foreign environment with computers and software they are not used to. In eye-tracking studies, participants must always keep a certain distance from the screen for data to be correctly collected. Also, the laboratory setting limits the recruitment of translators to a specific place and prevents researchers from obtaining a larger sample. In this study, it was hard to find the desired number of participants, especially those from marginalised groups such as non-binary individuals. This difficulty is common to research in the field (\cite{burtscher-2022-genderfairmt}) and reinforced the decision to conduct the translation experiment online. 

The analysis of translation and PE times as well as screen activity tracking delivered interesting results and insights into the gender-fair translation and PE process. However, asking participants to work with texts was fundamental to gaining deeper insights into their experience and/or point of view, which was discussed during retrospective interviews. The point of view of professional translators concerning the use of different GFL strategies was of utmost importance because of their cross- and meta-linguistic competence. Besides this, their knowledge is important because translation can change social relations and increase empowerment. A strict experimental setting was therefore beyond the scope of the present work and many variables were not controlled. As such, the results from this study cannot be generalised.

Furthermore, full experimental realism was not possible despite careful planning because (i) translators were required to work on a single screen in order for the whole translation and PE process to be recorded, and (ii) they received strict requirements concerning GFL, i.e. the use of a single strategy for each text. Limitations such as (i) are usual in similar research settings whereas (ii) pertains to the study objective, namely to investigate the translation and PE process of different strategies and so to allow translators to make informed decisions on how to translate GFL. 

The major limitation of the current study is, as explained, the lack of result generalisability. This limitation is common to experimental studies in TPR because funding constraints make it difficult to recruit enough participants, and the triangulation of different methods causes data explosion (\cite{o2009eye}). Even more so for a PhD project since it has to be carried out in a limited time by a researcher only with the support of their supervisors. As already mentioned, the number of study candidates was selected by reviewing critical literature in the field of TPR (e.g. \cite{o2009eye}) that illustrates how results obtained from a small number of 12 participants may have significant value although they lack generalisability. Furthermore, the sample of participants, though small, was relatively heterogeneous. Ideally, the group would be more homogeneous in order to control variables such as typing speed and work experience. Also, other languages should be investigated, possibly without English in the combination. 

Another limitation concerns the comparability of the texts and hence of the cognitive processes analysed. To ensure comparability of estimated translation and PE times, texts used in the experiment were controlled for length, i.e. each was about 150 words, and, in part, for difficulty by measuring their readability with the Flesch-Kincaid readability test (\cite{kincaid1975derivation}) as in \textcite{lopez2016can}. The test takes into account the number and length of words but ignores semantics. This issue was partially addressed by asking for translators' opinions on the source texts' difficulty. In an ideal scenario, the same text could have been translated or post-edited with different strategies by more language professionals. This would require a higher number of participants than available within the scope of this study. In addition, the assignment order was not changed among participants and the most challenging approach to GFL, i.e. neosystems, was applied in the third text. Longer translation and PE times could hence be caused not only by the specific strategy applied but also by participants' tiredness. Furthermore, the impact of participant familiarity with the texts and potential learning effects on assignments was not considered. Although the author did not specifically ask participants during the interview whether they were familiar with the TV series described in the assignments, most of them revealed that they were not. In addition, gender references in each text slightly differed in their number (nine, ten, and twelve respectively) and difficulty, (e.g. ``nanny'' in the third text for which German equivalents are usually of feminine gender). Some terms caused difficulties because there is no equivalent or they are not so common in German (e.g. ``recurring guest star''). In replicating the study, terminology from the field of cinema and TV series could be provided so as to distinguish more clearly where the effort lies, i.e. in translating a term or in using a gender-fair strategy. Finally, translation and PE screen activity tracking can only be compared if changes are not counted. Screen recordings are of interest though, because they show which gender-specific terms are particularly challenging for translators and possibly for MT debiasing.

Although the study focuses on the translation and PE of GFL beyond the binary, only one non-binary translator participated in the study. To ensure their inclusion in the study, non-binary individuals were directly contacted. While numerous people were found, they had no translation background or translation experience. They were therefore excluded from the study. This may be considered as a major limitation of the present work. However, its scope is not to prescribe the use of a specific GFL strategy. The main objectives are to further expand the field of queer TS and also understand the language professionals' perceptions and challenges in the use of GFL to better support them in their work. 

People who decided to participate in the study had usually positive attitudes toward GFL whilst others were interested in the topic because it is becoming increasingly relevant to their work life. This may have had an impact on the results. Effort may be higher in translators who do not or refuse to use GFL. This is quite plausible since the only translator in the study who did not use GFL in their work life was also one of the participants who faced the highest effort. 

The identification of gender phrases for the target text analysis was particularly challenging. While they can be easily identified in the source text, there is no one-to-one correspondence among languages and translation. This may also be one of the reasons why participants used rewording in the second and third assignments as well. During the annotation process, the author read the target texts and identified chunks containing GFL. In this work, the term \textit{gender phrase} was used as it refers to a group of words standing together as a conceptual unit. A similar concept, \textit{translation unit}, is often used in TPR. Since its conceptualisation is widely debated (e.g. \cite{carl2011gazing}), it was not adopted in the present study. The process of deciding which grouping of words constituted a gendered phrase is therefore subjective and other researchers could opt for a different categorisation. 

As discussed in Subsection \ref{sub:biasMT}, the present work was developed prior to the rise in interest on LLMs, hence the potential of generative AI for personalised, bias-mitigated translation was not addressed and this study only addresses limitations of ``classical'' NMT systems.

In a nutshell, though the present study has several limitations, it is a first step towards including non-binary language in (queer) TS. A similar setup can be used in future research to replicate this work, possibly with a larger sample of participants, and to investigate the gender-fair translation and PE process in other languages. Finally, other methodologies such as keylogging and eye-tracking could be used in order to gather more objective and precise data on cognitive processes. 

The minimal and not statistically significant difference in translation and PE times among assignments may also indicate that language professionals can adapt to different requirements such as the use of a specific language variant without an impact on their speed and productivity. Experienced professionals are expected to have this skill, but it would need to be tested explicitly by translating with and without gendering. However, the resulting cognitive effort seems relatively high and future research could and should clarify whether this has an impact on the general quality of the target texts. Other populations, e.g. students, might produce different results.

Finally, it would be interesting to test the acceptability of the strategies in end users, e.g. studies with readers, especially focusing on those who are most impacted by the use of GFL, i.e. non-binary individuals. Studies on the readability and comprehensibility of new gender-fair forms are needed as well. A major complication in this respect is the lack of a standard, widely adopted system. Language is constantly evolving and minoritised groups do not necessarily agree on a specific set of forms. Moreover, many people do not know what GFL is and why it is being used. Greater societal awareness is needed for gender-fair forms to be more widely understood and used.

